% Einheit 3 von 3 – Am Ratengraphen (bank/rekonstruktion-von-bestaenden/e3.jsonl, zusammenbau v0.1)
%% TODO Zeitmarke Einheit 3: Marken-Zeile nicht lesbar
%% TODO Zweigzeile Teil 1: was hier gelernt wird (Satz in Schülersprache aus dem Einheitstitel) – Einheit 3
\einheitenkopf[e3]{Einheit 3 von 3 · Am Ratengraphen}
\zweigzeile{Abitur GK}
\setcounter{aufgabe}{13}

%% TODO Titel: Ich-kann-Satz fehlt in der Bank (Platzhalter: Kettenname) – Nr. 14
%% TODO Anweisung: ein Satz über der ersten Teilaufgabe fehlt in der Bank – Nr. 14
\begin{aufgabe}{Am Ratengraphen}
%% TODO \rechenplatz{n}: Schrittzahl des Grundfalls fehlt in der Bank (nur bei fester Schreibform, 2.3 b)
\begin{teile}
\teil Die Abbildung zeigt die Zuflussrate $r$ eines Behälters in Litern je Minute. Lies $r(1)$ ab und deute den Wert. Um wie viel ändert sich der Inhalt von $t = 0$ bis $t = 4$?

\begin{ksys}[xmin=0,xmax=7,ymin=-3,ymax=5,xlabel=t,ylabel=r,ablesen] \funktionab{4-\x}{r}{0}{6} \end{ksys}
\teil Zwei Pflanzen wachsen mit $p(t) = 6t - t^2$ und $q(t) = 0{,}5t^2$ (cm je Woche, $t$ in Wochen), beide starten mit der Höhe 0. Berechne $∫_0^4 (p(t) - q(t))\,\mathrm{d}t$ und deute den Wert. \hfill (Abitur 2019 GK)
\teil Die Abbildung zeigt die Änderungsrate $r$ des Wassers in einem Becken in m³ je Stunde. Markiere den Zeitpunkt $T > 0$, zu dem wieder genauso viel im Becken ist wie zu Beginn, und begründe mit Flächen.

\begin{ksys}[xmin=0,xmax=8,ymin=-4,ymax=4,xlabel=t,ylabel=r,ablesen] \funktionab{3-\x}{r}{0}{7} \end{ksys}
\teil Ein Parkhaus ist von 7 bis 23 Uhr geöffnet, vorher und nachher steht kein Auto darin. $e$ ist die Einfahrrate, $a$ die Ausfahrrate (Autos je Stunde). Die Graphen schließen mit der $t$-Achse die Fläche I (nur unter $e$, links vom Schnittpunkt), II (unter beiden) und III (nur unter $a$, rechts vom Schnittpunkt) ein. Erläutere im Sachzusammenhang, dass I und III gleich groß sind. \hfill (Abitur 2026 LK)
\end{teile}
\end{aufgabe}

%% TODO Titel: Ich-kann-Satz fehlt in der Bank (Platzhalter: Kettenname) – Nr. 15
%% TODO Anweisung: ein Satz über der ersten Teilaufgabe fehlt in der Bank – Nr. 15
\begin{aufgabe}{Am Ratengraphen – Fehler finden, Begründen, Darstellungswechsel, Anwendung}
\begin{teile}
\teil Zwei gleich volle Becken erhalten die Zuflüsse $p(t) = 2t$ und $q(t) = 6$ (Liter je Minute). Tom sagt: „Bei $t = 3$ ist in beiden gleich viel, denn da schneiden sich die Graphen.“ Finde den Fehler und rechne richtig.
\teil Begründe: Ist ein Bestand zu zwei Zeitpunkten gleich, sind dazwischen die Flächen ober- und unterhalb der Achse gleich groß.
\teil Zeichne in das Koordinatensystem den Graphen einer Rate $r$ auf $[0; 6]$, bei der der Bestand zum Zeitpunkt $t = 6$ wieder so groß ist wie bei $t = 0$.

\begin{ksys}[xmin=0,xmax=7,ymin=-4,ymax=4,xlabel=t,ylabel=r] \end{ksys}
\teil Ein Speicher enthält zu Beginn 200\,m³, er erhält konstant 50\,m³ je Stunde und gibt $a(t) = 8t$ (m³ je Stunde) ab, $0 \le t \le 10$. Wie viel ist nach zehn Stunden im Speicher, und wann war am meisten darin?
\end{teile}
\end{aufgabe}
